\documentclass[%
 aip,
% jmp,
% bmf,
% sd,
% rsi,
 amsmath,amssymb,
preprint,%
% reprint,%
%author-year,%
%author-numerical,
]{revtex4-2}

\usepackage{graphicx}
\usepackage{dcolumn}
\usepackage{bm}
%\linenumbers\relax % Commence numbering lines

\usepackage[utf8]{inputenc}
\usepackage[T1]{fontenc}
\usepackage{mathptmx}
\usepackage{etoolbox}

%% Apr 2021: AIP requests that the corresponding 
%% email to be moved after the affiliations
\makeatletter
\def\@email#1#2{%
 \endgroup
 \patchcmd{\titleblock@produce}
  {\frontmatter@RRAPformat}
  {\frontmatter@RRAPformat{\produce@RRAP{*#1\href{mailto:#2}{#2}}}\frontmatter@RRAPformat}
  {}{}
}%
\makeatother

\begin{document}


\title{Molecular responses in a snap: the \texttt{std2} package}

\author{Pierre Beaujean}
 \affiliation{
Theoretical Chemistry Laboratory, Unit of Theoretical and Structural Physical Chemistry, Namur Institute of Structured Matter (NISM), University of Namur, Rue de Bruxelles, 61, B-5000 Namur (Belgium)
 }
 
\author{Stephan Grimme}%
\affiliation{ 
	Mulliken Center for Theoretical Chemistry, Institut für Physikalische und Theoretische Chemie, Beringstr. 4, G-53115 Bonn (Germany)
}

\author{Marc de Wergifosse}
\email{marc.dewergifosse@uclouvain.be}
\affiliation{
	Theoretical Chemistry Group, Molecular Chemistry, Materials and Catalysis Division (MOST), Institute of Condensed Matter and Nanosciences, Université Catholique de Louvain, Louvain-la-Neuve B-1348 (Belgium)
}

\date{\today}

\begin{abstract}
An abstract.
\end{abstract}

\maketitle

\section{Introduction}

We need sQC for large system and/or snapshots of smaller system + history of the different flavors of sQC, using Refs.~\citenum{dewergifossePerspectiveSimplifiedQuantum2021}, \citenum{loffelsenderSimplifiedQuantumChemistry2024}, and \citenum{dewergifosseEXactIntegralSimplified2024}.

\section{Theory}

Using Ref.~\citenum{loffelsenderSimplifiedQuantumChemistry2024}:

\begin{itemize}
	\item DFT ;
	\item TD-DFT, then simplified TD-DFT, from XsTD to sTD ;
	\item NTOs and response orbitals ;
	\item Implementation details (Fortran, libcint, etc).
\end{itemize}

\section{Illustrations}

\subsection{Obtaining the ground state}

... And interfacing with other programs.

\subsection{Excitations and 1PA spectra}

\texttt{stda} was originally designed for UV and ECD spectra \cite{grimmeSimplifiedTammDancoffDensity2013,bannwarthSimplifiedTimedependentDensity2014,grimmeUltrafastComputationElectronic2016,loffelsenderAllAtomQuantumMechanical2025} (and then with XsTD \cite{dewergifosseEXactIntegralSimplified2024,dewergifosseComputingExcitedStates2024,maraldiEvaluatingPerformanceExact2025}), but was later extended to spin-flip \cite{dewergifosseSimplifiedSpinFlipTimeDependent2019,dewergifosseAreFullyConjugated2019}, XAS [?], and ESA \cite{dewergifosseNonlinearresponsePropertiesSimplified2019}.

\subsection{Linear responses}

Optical rotations \cite{dewergifosseSimplifiedTimedependentDensity2020} and alpha.

\subsection{Quadratic responses}

beta \cite{dewergifosseNonlinearresponsePropertiesSimplified2018,seibertDynamicStructuralEffects2020,beaujeanAllAtomQuantumMechanical2021,lescosSelfassemblingStructureNonlinear2021,dewergifosseEXactIntegralSimplified2024} and 2PA spectra \cite{dewergifosseUltrafastEvaluationTwoPhoton2022,dewergifosseEXactIntegralSimplified2024,dewergifosseComputingExcitedStates2024,loffelsenderAllAtomQuantumMechanical2025,maraldiEvaluatingPerformanceExact2025}

\subsection{Post-treatments}

NTOs and response orbitals \cite{dewergifosseUnifiedStrategyChemically2020}.

\section{Conclusions and perspectives}

\begin{itemize}
	\item ``\texttt{std2} provides approximations that retain most of the physics" [citation needed]
	\item What about AI?
\end{itemize}

\begin{acknowledgments}
	We acknowledge past and future devs.
\end{acknowledgments}

\section*{Data Availability Statement}

Data sharing is not applicable to this article as no new data were created or analyzed in this study.

\bibliography{biblio}

\end{document}
